\hyt{vreledikyzpupku}
\song{Vřelé díky z pupku} \interpret{zuzananavarova}{Zuzana Navarová}

\vers{1}{
A \chord{E\7}šajnym dank\footnote{z němčiny \uv{schönen dank} = mnohokrát děkuji} in \chord{Am}pupik, aj \chord{G}aj aj mazl-\chord{C}tov\footnote{\uv{hodně štěstí}, doslova \uv{dobrá hvězda} – z hebrejského \textit{mazal} (\uv{souhvězdí}, \uv{osud}) a \textit{tov} (\uv{dobrý}). Používá se při svatbách, narozeních, úspěších apod.},\\
a \chord{E\7}šajnym dank in \chord{Am}pupik, líp ti \chord{H\7}nezacinká \chord{E\7}žádný kov.\\
A \chord{E\7}šajnym dank in \chord{Am}pupik, Bože \chord{G}chraň, ó Ado\chord{C}naj\footnote{Hebrejské slovo \uv{můj Pán}. V židovské tradici se při čtení biblického textu vyslovuje místo Božího jména JHVH.},\\
a \chord{E\7}šajnym dank in \chord{Am}pupik, jak by \chord{E\7}chasene\footnote{Jidiš výslovnost hebrejského \textit{chatuna} (\uv{svatba}). Ve východoevropské židovské kultuře je chasene velká rodinná událost s hudbou, tancem a hostinou.}, tak \chord{Am}hraj, hraj, hraj.
}

\refrain{
Hú \chord{Dm}ha, to je mi hudba sfér, hú \chord{Am}ha, to se raduje sám Ahasver\footnote{Ahasver je postava \uv{věčného Žida}, který kvůli svému provinění musí až do dne posledního soudu bloudit po zemi. Podle legendy Ahasver udeřil Ježíše Krista při výslechu u Kaifáše. Podle jiné verze se mu vysmíval, když nesl kříž, slovy: \uv{Tak ty tvrdíš, že se vrátíš.} A Ježíš odpověděl: \uv{Ano, a ty tu na mě počkáš.} A od té doby musí Ahasver bloudit po zemi a čekat. Objevuje se v dílech Byrona či Vrchlického, též Apollinaira -- Pražský chodec (Isaac Lakedem), nebo v novele Krysař od Viktora Dyka. V přeneseném významu jméno označuje nespokojeného, věčně se toulajícího člověka, zarostlého a nuzného vzhledu, který nenachází uspokojení a uplatnění.},\\
hú \chord{E\7}ha, s ním hnedle cha cha cha dělá \chord{E\7}celá mišpa\chord{Am}cha\footnote{= hebrejsky \uv{rodina}}.
}

\vers{2}{
A šajnym dank in pupik, aj aj aj mechaje\footnote{Z hebrejského \textit{mechaje} (oživující). V jidiš může znamenat něco jako \uv{to je paráda!}},\\
a šajnym dank in pupik, kdo tu není klezmer\footnote{Původně z hebrejských slov \textit{kli} (\uv{nástroj}) + \textit{zemer} (\uv{zpěv}, \uv{melodie}), tedy \uv{hudební nástroj}. Později označení pro muzikanta a celý hudební styl.} nehraje.\\
A šajnym dank in pupik, to hlas můj jako zvon,\\
Bůh přisám na sám pupik, a tak alejchem šalom\footnote{\textit{šalom alejchem} je tradiční židovský pozdrav. Doslova \uv{mír vám}. Odpověď zní \textit{alejchem šalom} (\uv{vám mír}).}.
} \refsm{}

\bridge{
Oj oj oj, kdo ví, kde nota gé? Oj oj oj, já nebich mešuge\footnote{= blázen, šílený}.
} \refsm{} $\times 2$ \bridge{}
\ns

\vers{3}{
A kdo \chord{E\7}nerozuměl \chord{Am}jidiš, a kdo \chord{G}nikdy nejed' \chord{C}šoulet\footnote{Hustý pokrm z fazolí, ječmene, masa a brambor, připravovaný přes noc před šabatem, aby byl v sobotu hotový bez vaření.},\\
kdo se \chord{E\7}divil jako \chord{Am}děcko, bude \chord{G}po něm bulvy \chord{C}koulet,\\
i když \chord{E\7}nedojí to \chord{Am}všecko.
} \refsm{} $\times 2$

\newpage
